\section{Results}
\label{sec:results}

This section reports fixed-policy results for Qwen2.5 and Qwen3 families
across two consumer-GPU setups: desktop RTX4090 and RTX5090-laptop.
We focus on paired speedup $S$, acceptance dynamics ($\alpha$, $B_{\text{eff}}$),
and token-timing metrics (time to first token, TTFT; time per output token,
TPOT), then summarize portability implications.

\begin{table*}[t]
\centering
\caption{Qwen2.5 fixed-policy results on RTX4090 (target: 3B, draft: 0.5B) against regime-matched autoregressive baselines.
Speedup is computed as a ratio of means against the corresponding baseline.}
\label{tab:headline-4090-qwen25}
\small
\begin{tabular}{l c c c c c c}
\hline
\textbf{Config} & $S$ & $\alpha$ & $B_{\text{eff}}$ & $T_{\text{mean}}$ (s) & TTFT (ms) & TPOT (ms) \\
\hline
0.5B, $k{=}4$, det   & \textbf{3.0674} & 0.4212 & 1.66 & 3.7210 & 102.35 & 29.06 \\
0.5B, $k{=}4$, stoch & 2.9394 & 0.4080 & 1.61 & 3.9631 & 105.26 & 30.89 \\
0.5B, $k{=}8$, det   & 2.4542 & 0.3515 & 2.72 & 4.6508 & 162.76 & 36.73 \\
0.5B, $k{=}8$, stoch & 2.4063 & 0.3444 & 2.67 & 4.8411 & 165.95 & 38.04 \\
0.5B, $k{=}16$, det  & 1.8050 & 0.2539 & 3.76 & 6.3234 & 277.43 & 49.34 \\
0.5B, $k{=}16$, stoch& 1.8185 & 0.2516 & 3.73 & 6.4059 & 278.85 & 49.81 \\
\hline
\end{tabular}
\end{table*}

\begin{table*}[t]
\centering
\caption{Qwen3 fixed-policy results on RTX4090 against regime-matched autoregressive baselines.
Speedup is computed as a ratio of means against the corresponding baseline.}
\label{tab:headline-4090}
\small
\begin{tabular}{l c c c c c c}
\hline
\textbf{Config} & $S$ & $\alpha$ & $B_{\text{eff}}$ & $T_{\text{mean}}$ (s) & TTFT (ms) & TPOT (ms) \\
\hline
0.6B, $k{=}4$, det   & \textbf{2.8548} & 0.3362 & 1.33 & 4.8881 & 132.57 & 41.23 \\
0.6B, $k{=}4$, stoch & 2.4798 & 0.2672 & 1.06 & 5.6196 & 134.54 & 45.51 \\
0.6B, $k{=}8$, det   & 1.9648 & 0.2378 & 1.86 & 7.1024 & 204.99 & 58.76 \\
0.6B, $k{=}8$, stoch & 1.6869 & 0.1916 & 1.50 & 8.2612 & 207.82 & 64.67 \\
0.6B, $k{=}16$, det  & 1.1091 & 0.1292 & 1.96 & 12.5817 & 349.58 & 94.24 \\
0.6B, $k{=}16$, stoch& 0.9719 & 0.1082 & 1.64 & 14.3390 & 354.82 & 106.03 \\
\hline
\end{tabular}
\end{table*}

\subsection{RQ1: Does speculative decoding accelerate inference?}

On RTX4090, Qwen2.5 and Qwen3 both show clear acceleration in most fixed-policy
settings. For Qwen2.5 (Table~\ref{tab:headline-4090-qwen25}), all six reported
points exceed baseline, with best speedup at 0.5B, $k{=}4$, deterministic
($S{=}3.0674$). For Qwen3 (Table~\ref{tab:headline-4090}), five of six points
exceed baseline, with best speedup at 0.6B, $k{=}4$, deterministic
($S{=}2.8548$). The one below-baseline case is Qwen3 high-$k$ stochastic
($k{=}16$, $S{=}0.9719$), indicating that large proposal length can remove net
latency gains.

\subsection{RQ2: How do $k$ and acceptance affect net latency?}

Both families retain the same systems trend: larger $k$ increases accepted
block size but degrades end-to-end speed. On RTX4090, Qwen2.5 mean speedup by
$k$ is $\bar{S}_{k=4}=2.9164$, $\bar{S}_{k=8}=2.4303$,
$\bar{S}_{k=16}=1.8118$, while Qwen3 shows
$\bar{S}_{k=4}=2.6673$, $\bar{S}_{k=8}=1.8258$,
$\bar{S}_{k=16}=1.0405$. In both tables, $B_{\text{eff}}$ rises with $k$
while speedup falls, indicating an overhead inflection where additional
drafting and verification cost dominates block-size benefits.

\subsection{RQ3: Deterministic versus stochastic regime}

For Qwen3, deterministic decoding is faster than stochastic at every matched
$k$: $\Delta S_{k=4}{=}0.3750$, $\Delta S_{k=8}{=}0.2779$,
$\Delta S_{k=16}{=}0.1372$ (deterministic minus stochastic), with mean gap
0.2633 speedup units (Table~\ref{tab:headline-4090}).
For Qwen2.5, deterministic is also faster at $k{=}4$ and $k{=}8$
(Table~\ref{tab:headline-4090-qwen25}), with a small reversal at $k{=}16$.
Across both families, deterministic generally provides a stronger runtime
profile at low-to-mid $k$.

\subsection{RQ4: Cross-hardware portability (RTX4090 vs RTX5090-laptop)}

\begin{table}[t]
\centering
\caption{Qwen3 fixed-policy results on RTX5090-laptop.
Ordering trends are preserved (best at low $k$, deterministic $>$ stochastic),
but absolute speedups are below 1 under the tested runtime stack.}
\label{tab:headline-5090}
\small
\begin{tabular}{l c c c}
\hline
\textbf{Config} & $S$ & $\alpha$ & $B_{\text{eff}}$ \\
\hline
0.6B, $k{=}4$, det   & 0.6743 & 0.3353 & 1.32 \\
0.6B, $k{=}4$, stoch & 0.5924 & 0.2845 & 1.06 \\
0.6B, $k{=}8$, det   & 0.4613 & 0.2538 & 1.85 \\
0.6B, $k{=}8$, stoch & 0.3995 & 0.2273 & 1.50 \\
0.6B, $k{=}16$, det  & 0.2588 & 0.1803 & 1.96 \\
0.6B, $k{=}16$, stoch& 0.2222 & 0.1543 & 1.64 \\
\hline
\end{tabular}
\end{table}

The ordering pattern transfers across devices (low-$k$ best,
deterministic preferred, and monotonic decline with larger $k$),
but absolute gains do not.
Under the tested RTX5090-laptop software/runtime stack, all reported
fixed-policy points are slower than autoregressive baseline.
This indicates that
conclusions are more portable than absolute speedup magnitudes.

\subsection{Variance, TTFT/TPOT, and statistical testing detail}

To provide stronger statistical detail, we include seed-level summaries from
repeated fixed-$k$ runs.
Across seeds, deterministic speedup means are
$2.2876\pm0.0242$ ($k{=}4$), $2.2712\pm0.0295$ ($k{=}8$), and
$2.0511\pm0.0273$ ($k{=}16$); stochastic means are
$1.8910\pm0.0220$, $1.9055\pm0.0257$, and $1.6698\pm0.0133$.
TTFT is stable (about 50--52 ms), while TPOT and 99th-percentile (P99)
latency rise with larger
$k$, consistent with the overhead interpretation.

We also report paired tests with Holm-Bonferroni correction for adaptive-$k$-versus-fixed
comparisons in the adaptive-$k$ run set, separating statistically
significant token-timing effects from non-significant quality differences in
that comparison set.

\subsection{Summary of practical operating point}

For the Qwen3 evidence, the most reliable deployment recommendation on
RTX4090 remains deterministic $k{=}4$ as the default fast path.
Cross-device validation supports this policy ranking but warns that
hardware/runtime portability is limited for absolute throughput gains.

\subsection{Cross-family consistency (Qwen2.5 vs Qwen3)}

To address model-family applicability, we compare Qwen2.5 RTX4090
results with the Qwen3 results.
The absolute speedup scale differs across runs, but the core qualitative
patterns are consistent across both families on RTX4090:
(i) best operating point at low $k$ (especially $k{=}4$),
(ii) worsening speedup as $k$ increases, and
(iii) deterministic preference in most matched settings.

Concretely, the Qwen2.5 table reports best fixed speedup at
0.5B, $k{=}4$, deterministic ($S{=}3.0674$), with mean-$k$ trend
$\bar{S}_{k=4}=2.9164 > \bar{S}_{k=8}=2.4303 > \bar{S}_{k=16}=1.8118$.
In the Qwen3 RTX4090 run family, we observe the same ordering with
different magnitudes ($\bar{S}_{k=4}=2.6673 > \bar{S}_{k=8}=1.8258 >
\bar{S}_{k=16}=1.0405$).

These aligned trend-level findings support a bounded claim:
operating-policy ranking appears more portable across these two model
families than absolute throughput values, which remain implementation- and
hardware-path dependent.
